% SPDX-License-Identifier: MIT
\chapter{Praxis: eine administrative Bedienoberflaeche}\label{chap:practice}
\ChapterSummary{Sie planen eine kleine Administrationstool-
Oberflaeche nach den Prinzipien des SASD UI Toolkit:
ein Datenmodell für den Zustand, eine kontrollierte Auswahl,
Commands für Änderungen und eine Darstellung, die
den Benutzer nicht über ihre Berechtigungen taeuscht.}

\section{Beispielaufgabe und Grundsaetze}
Angenommen, eine Anwendung soll eine Liste
von Diensten anzeigen, deren Status lesen
und einen ausgewaehlten Dienst neu laden.
Für SASD Gatehold waere dieselbe
Struktur für Firewallregeln oder
Netzwerkschnittstellen geeignet.
Wir unterscheiden vier Ebenen:
\begin{enumerate}
\item Die \term{Domaene} kennt den wirklichen
Zustand des Systems und validiert Aktionen.
\item Das \term{Model} liefert einen
lesbaren Ausschnitt der Domaenendaten.
\item Die \term{View} verwaltet Darstellung
und aktuelle Auswahl.
\item Ein \term{Command} beschreibt
eine anforderungsgesteuerte Aktion.
\end{enumerate}
Die UI darf anzeigen, dass ein Command
nicht verfügbar ist. Das ist eine
wichtige Rückmeldung, aber keine
Authentisierung oder Autorisierung.

\section{Schritt eins: nur lesen}
Beginnen Sie mit einem \term{StringTableModel}
aus \texttt{Dienst} und \texttt{Status}.
Fuellung und Erneuerung des Modells
geschehen getrennt vom Renderer.
Eine beispielhafte Quelle könnte einen
Snapshot nach erfolgreicher Statusabfrage
zurueckgeben. Ein \term{TableView}
zeigt die ersten sichtbaren Zeilen;
das \term{TableSelectionModel}
identifiziert die aktuell gewählt Zelle.

\begin{lstlisting}[language=SASDCpp]
StringTableModel services(
    {"Dienst", "Status"},
    {{"dns", "aktiv"}, {"proxy", "aktiv"}});
TableSelectionModel selection;
selection.setModel(&services);

TableView view;
view.setModel(&services);
view.setSelectionModel(&selection);
view.setViewport(0, 10, 0, 2);
\end{lstlisting}
In einem realen Fenster muss die View in
einen Container aufgenommen werden.
Dieser Ausschnitt demonstriert den
Datenvertrag unabhaengig vom Host.
Er ersetzt weder den Eventloop
noch die Terminal-Integration.

\section{Schritt zwei: Aktion mit Guard}
Ein \term{Command} namens
\emph{Dienst neu laden} soll nur bei
einem passenden Datensatz aktiviert werden.
Dabei müssen mindestens folgende
Fragen fachlich beantwortet sein:
Existiert der Datensatz noch?
Ist der Status frisch?
Ist die Aktion für diesen Dienst erlaubt?
Hat der Ausführende die Berechtigung?
Ist eine Bestätigung erforderlich?

Die UI kann \texttt{setEnabled(false)}
setzen, bis alle sichtbaren Vorbedingungen
gegeben sind. Unmittelbar vor der
tatsächlichen Änderung müssen die
sicherheitsrelevanten Bedingungen im
zustandstragenden Dienst \emph{erneut}
geprueft werden. Ein alter Viewport
oder ein verzoegertes Tastaturereignis
darf nicht zum unkontrollierten
Neustart eines anderen Dienstes fuehren.

\section{Schritt drei: Rückmeldung und Fehlerbild}
Aus einer administrativen Aktion
müssen für Benutzer erkennbare
Zustände hervorgehen: \emph{angefordert},
\emph{in Arbeit}, \emph{erfolgreich},
\emph{fehlgeschlagen} und gegebenenfalls
\emph{zurueckgenommen}.
Ein \term{Label} kann den aktuellen
Status anzeigen; eine
\term{ListView} könnte eine
kleine Ereignisliste darstellen.
Die echte Protokollierung und
Auditierung verbleiben ausserhalb
der UI-Bibliothek.

\section{Schritt vier: robuste Eingabe}
Testen Sie nicht nur den Gluecksfall.
Was passiert beim Resize während
eines offenen Popupmenues?
Was bei einem Model-Reset genau nach
der Zeilenauswahl?
Was bei einer Aktion, die das aktuelle
Fenster schliesst?
Was bei einem Terminal ohne Maus?
Was passiert nach Verlust des
Terminal-/SSH-Kanals?
Eine nützliche Verwaltungsoberflaeche
erhält für solche Fragen explizite
Antworten und automatisierte Tests.

\section{Reale Projektreihenfolge}
Für ein erstes wartbares Werkzeug ist
diese Reihenfolge sinnvoll:
\begin{enumerate}
\item \textbf{Lesepfad:} Daten korrekt
laden und ohne Veränderung anzeigen.
\item \textbf{Fokuspfad:} Vollständige
Tastaturbedienung nachweisen.
\item \textbf{Selektion:} Datensatz
und Auswahl sichtbar abgleichen.
\item \textbf{Commands:} Erst lokale
Simulation, dann echte, abgesicherte
Systemoperation.
\item \textbf{Fehlerpfad:} Status,
Auditlog und Wiederherstellbarkeit.
\item \textbf{Zweiter Backend-Pfad:}
Wenn erforderlich dieselbe Semantik
mit Rendered testen.
\end{enumerate}

\PracticalTip{Nehmen Sie zum Start eine
\emph{read-only} Bildschirmseite,
beispielsweise eine Dienstliste.
Die Kombination aus TableModel,
SelectionModel und einem deaktivierten
Command reicht aus, um den
gesamten Daten- und Interaktionsweg
ohne Risiko zu proben.}
